\begin{theorem}[Bool CantorSpace endpoint route]
\label{thm:bool-cantor-space-endpoint-route}
Every displayed $\BoolUp$ endpoint row consumed by a $\CantorSpaceUp$ finite prefix cylinder remains a $\BoolUp$ two-endpoint $\NameCert$ row after $\StreamNameUp$ and $\ListUp$ prefix transport. The route exports only the empty and one-step endpoint rows, their history classifier, the transported prefix position, and the local package ledger. It does not introduce a host Boolean, quotient stream equality, branch choice, or compactness claim.
\end{theorem}
\begin{proof}
Start from the public endpoint surface of $\BoolUp$, where \autoref{def:bool-public-endpoint-constructors} displays exactly the two rows $\emp$ and $\Eone(\emp)$. The case readback theorem \autoref{thm:bool-public-constructor-case-readback} fixes every public endpoint read to one of these rows, with agreement measured by $\msame$ and the endpoint classifier rather than by a host Boolean predicate.

The Cantor prefix consumer in \autoref{ch:concrete-instances-cantorspace-namecert} reads Boolean observations only through finite prefix cylinders. Its StreamName schedule is the window discipline of \autoref{ch:concrete-instances-streamname-namecert}, and its finite prefix spine is a list-indexed transport of displayed endpoint rows. Transport along that schedule changes the inspected position and the finite prefix ledger, not the two endpoint constructors supplied by $\BoolUp$.

Thus the transported row remains a $\BoolUp$ endpoint row after the CantorSpace prefix and StreamName/List bookkeeping have been attached. The no-hidden-leakage theorem \autoref{thm:bool-public-surface-no-hidden-leakage} excludes host Boolean structure, while the CantorSpace chapter supplies only the finite-prefix consumer route and not quotient streams, branch choices, or compactness.
\end{proof}
